\begin{ouverture}[calcul]{2}{Fractions et nombres premiers}{Partie 1 · Nombres et calculs}
{Nombres rationnels · Inverse · Simplification · Partages}
\objectif{Calculer avec toutes les opérations sur les fractions.}
\objectif{Comparer des rationnels et simplifier avec les nombres premiers.}
\objectif{Résoudre des problèmes de partage et de proportion.}
\prerequis{Multiples, diviseurs et règles des signes.}
\end{ouverture}
\pagestyle{otMain}\essentiel\fontsize{11}{15}\selectfont
\phantomsection\label{cours-2-1}
\essSection{1}{Écrire et comparer}
\begin{definition}Un nombre rationnel peut s’écrire $a/b$ où $a$ et $b$ sont des entiers et $b\ne0$. Une fraction représente un nombre, pas deux nombres séparés. Multiplier ou diviser son numérateur et son dénominateur par un même nombre non nul conserve sa valeur.\end{definition}
\begin{propriete}Pour comparer des fractions, on peut les écrire avec un même dénominateur positif. Pour additionner ou soustraire, utiliser ce dénominateur commun puis opérer sur les numérateurs : $\frac ab+\frac cb=\frac{a+c}b$.\end{propriete}
\begin{definition}Un nombre premier a exactement deux diviseurs positifs : 1 et lui-même. Ainsi 1 n’est pas premier. Décomposer un entier en facteurs premiers aide à repérer les facteurs communs. Une fraction est irréductible quand les deux entiers n’ont plus de diviseur commun supérieur à 1.\end{definition}
\begin{exemple}$84=2^2\times3\times7$ et $90=2\times3^2\times5$. Donc $\frac{84}{90}=\frac{14}{15}$. Pour comparer $-2/3$ et $-3/4$, écrire $-8/12>-9/12$.\end{exemple}
\begin{savoirfaire}[Retrouver les premiers inférieurs à 100]Dans la liste de 2 à 99, barrer les multiples de 2, 3, 5 et 7, sauf ces quatre nombres eux-mêmes. Les nombres restants sont premiers :\par
{\small $2,3,5,7,11,13,17,19,23,29,31,37,41,$\par
$43,47,53,59,61,67,71,73,79,83,89,97.$}\par
Ce crible suffit ici : un entier composé inférieur à 100 possède un facteur premier inférieur à 10.\end{savoirfaire}
\clearpage
\phantomsection\label{cours-2-2}
\essSection{2}{Multiplier et diviser}
\begin{propriete}Pour multiplier : $\frac ab\times\frac cd=\frac{ac}{bd}$, avec $b,d\ne0$. Simplifier des facteurs avant de multiplier peut alléger le calcul. L’inverse d’un nombre non nul $a/b$ est $b/a$ : leur produit vaut 1.\end{propriete}
\begin{propriete}Diviser par un nombre non nul revient à multiplier par son inverse : $\frac ab:\frac cd=\frac ab\times\frac dc$, avec $b,c,d\ne0$. Zéro n’a pas d’inverse.\end{propriete}
\begin{exemple}$\frac{-3}{5}:\frac9{10}=\frac{-3}{5}\times\frac{10}{9}=-\frac23$. En revanche $\frac12+\frac13=\frac36+\frac26=\frac56$, et non $\frac25$.\end{exemple}
\begin{savoirfaire}[Une fraction d’une quantité]« Les $2/5$ de 30 » signifie $(2/5)\times30=12$. Si une fraction porte sur un reste, déterminer d’abord ce reste. Conserver les valeurs exactes jusqu’au résultat final et vérifier que les parts forment bien le tout.\end{savoirfaire}

\clearpage\rubrique{Découverte}
\begin{blocLivre}{ex-2-1}\begin{exercice}[1]{Des fractions égales}
Compléter $\frac{-3}{4}=\frac{\,?\,}{20}$ et simplifier $\frac{42}{56}$.
\end{exercice}\end{blocLivre}
\begin{blocLivre}{ex-2-2}\begin{exercice}[1]{Comparer deux négatifs}
Comparer $-5/6$ et $-7/9$.
\end{exercice}\end{blocLivre}
\begin{blocLivre}{ex-2-3}\begin{exercice}[1]{Additionner}
Calculer $\frac34-\frac56$.
\end{exercice}\end{blocLivre}
\begin{blocLivre}{ex-2-4}\begin{exercice}[1]{Multiplier}
Calculer $\frac{-4}{7}\times\frac{21}{10}$.
\end{exercice}\end{blocLivre}
\begin{blocLivre}{ex-2-5}\begin{exercice}[1]{Trouver un inverse}
Donner l’inverse de $-5/3$, puis calculer $\frac23:\frac{-5}{3}$.
\end{exercice}\end{blocLivre}

\rubrique{Entraînement}
\begin{blocLivre}{ex-2-6}\begin{exercice}[2]{Décomposer}
Décomposer 180 et 252 en facteurs premiers, puis simplifier $180/252$.
\label{fig-ex-2-6}\figureGeo{facteurs-180}
\end{exercice}\end{blocLivre}
\begin{blocLivre}{ex-2-7}\begin{exercice}[2]{Les nombres premiers}
Lister tous les nombres premiers strictement compris entre 70 et 100. Expliquer pourquoi 91 n’en fait pas partie.
\end{exercice}\end{blocLivre}
\begin{blocLivre}{ex-2-8}\begin{exercice}[2]{Un calcul enchaîné}
Calculer $\frac56-\frac23\times\frac94$.
\end{exercice}\end{blocLivre}
\begin{blocLivre}{ex-2-9}\begin{exercice}[2]{Diviser une somme}
Calculer $\left(\frac34-\frac12\right):\frac58$.
\end{exercice}\end{blocLivre}
\begin{blocLivre}{ex-2-10}\begin{exercice}[2]{Encadrer}
Encadrer $-17/5$ entre deux entiers consécutifs et ranger $-17/5$, $-10/3$, $-7/2$.
\end{exercice}\end{blocLivre}

\rubrique{Problèmes}
\begin{blocLivre}{ex-2-11}\begin{exercice}[3]{Lire une portion du reste}
Un réservoir contient 60 L. On utilise $2/5$ de son contenu, puis $1/3$ de ce qui reste. Combien de litres reste-t-il ?
\end{exercice}\end{blocLivre}
\begin{blocLivre}{ex-2-12}\begin{exercice}[3]{Des sachets identiques}
On répartit sans reste 84 perles bleues et 126 vertes dans des sachets identiques. Donner tous les nombres de sachets possibles et la composition quand on en fait le plus.
\end{exercice}\end{blocLivre}
\begin{blocLivre}{ex-2-13}\begin{exercice}[3]{Retrouver une quantité}
Les $7/12$ d’un budget valent 175 euros. Quel est le budget total et quelle somme n’a pas encore été utilisée ?
\end{exercice}\end{blocLivre}
\begin{blocLivre}{ex-2-14}\begin{exercice}[3]{Une affirmation à discuter}
Tom affirme : « Diviser un nombre positif par une fraction le rend toujours plus petit. » Donner un contre-exemple et préciser quand le quotient devient plus grand.
\end{exercice}\end{blocLivre}
\begin{blocLivre}{ex-2-15}\begin{exercice}[3]{Trois parts successives}
Une somme est partagée : Lina reçoit $1/4$ du total, Sami $2/3$ du reste et Zoé les 36 euros restants. Trouver le total et les deux premières parts.
\end{exercice}\end{blocLivre}

\clearpage
